\section{Panel comparisons when the intervention, baseline and population differ}
\label{sec:exec28-panel-designs}

A branch rollout is not a single causal contrast until its comparison population and information clock are fixed.
Suppose a bank introduces a servicing platform across regions while a fee concession arrives on a different calendar.
The sponsor wants the platform's incremental contribution to customer cash flow. An event-study graph may instead describe
a change relative to a moving baseline, a combined platform-and-concession effect, or an average over a changing set of regions.
These are different questions even when every output column is called an ATT. The design therefore precedes the choice of R,
Stata or a language-neutral adapter. Three primary treatments make the distinction concrete: event-study normalization,
multiple events, and balanced stacked comparisons.~\citep{exec28_panel_roth,exec28_panel_tsai,exec28_panel_wing}

\subsection{What the event-study vertical axis measures}

Let treatment start at calendar time $t=1$, let $r=t-1$ denote event time, and let
$\Delta_t=E[Y_t\mid D=1]-E[Y_t\mid D=0]$ be the observed treated-minus-control mean gap.
In the common-adoption setting, a dynamic two-way fixed-effects comparison with period zero as its reference has
$\beta_r=\Delta_{r+1}-\Delta_0$. Roth's selected version shows why this interpretation cannot automatically be applied
to every pre-treatment coefficient from another estimator. A varying-base construction uses $\Delta_{r+1}-\Delta_r$
before treatment and $\Delta_{r+1}-\Delta_0$ afterwards. An imputation construction in his example uses the earliest
pre-period for pre-treatment coefficients but the mean pre-treatment outcome for post-treatment coefficients.
Different baselines create different graphical shapes even without a treatment effect.~\citep{exec28_panel_roth}

Consider an explicitly synthetic banking example, measured in dollars per customer per period. Treatment has zero effect,
but the untreated group gap is $\Delta_t=50t$ for $t=-4,\ldots,2$. This violates parallel trends deliberately. The last
pre-period is zero, the earliest gap is $-200$, and the mean pre-period gap is $-100$. At calendar times $-2,-1,0,1,2$,
the fixed-base contrasts are $-100,-50,0,50,100$. The varying-base contrasts are $50,50,50,50,100$ where the adjacent
pre-period exists. The earliest-base/mean-base imputation contrasts are $100,150,200,150,200$.
The apparent downward jump at adoption is a change in reference, not a negative platform effect. All these values
are independently recomputed from exact fractions.~\citep{exec28_panel_reference}

The practical remedy is to publish the comparison formula with the graph. In Roth's common-adoption example, a universal
pre-period base makes the Callaway--Sant'Anna construction numerically match the fixed-base comparison. Using the mean
pre-period on both sides gives $\Delta_t+100$ here, a constant vertical shift of the fixed-base series.
That pre-period series averages to zero by construction, so its average cannot independently establish parallel trends.
Roth also distinguishes in-sample pre-treatment residuals from post-treatment prediction errors: in his common-adoption
analysis the former can shrink toward zero. A small pre-period residual is not a general validation of the missing
counterfactual.~\citep{exec28_panel_roth}

This does not invalidate the recent estimators when their identifying assumptions hold. It addresses their diagnostic
behavior when those assumptions fail. The selected Roth paper records installations on 31 October 2025 and distinguishes
an earlier from an updated multiple-groups implementation. The separately pinned January 2024 example script has unpinned
package imports and cannot establish that later environment. A proposal should retain the diagnostic construction and
environment separately, and require a dependency lock before a future native replay. For the bank, the useful output is
the cohort path, its baseline and uncertainty, not merely an attractive common plot.~\citep{exec28_panel_roth,exec28_panel_software}

\subsection{A second event changes both the controls and the target}

Write $G^1$ for platform adoption and $G^2$ for concession adoption. Following the selected two-event framework, each
treatment is absorbing; reversal is outside this design. Let $Y_t(g^1,g^2)$ denote the potential outcome under both adoption
dates, with zero in an argument meaning no adoption of that intervention. For the joint cohort $G=(g^1,g^2)$ define
\begin{align*}
 C_G(t)&=E[Y_t(g^1,g^2)-Y_t(0,0)\mid G],\\
 A_G(t)&=E[Y_t(g^1,0)-Y_t(0,0)\mid G].
\end{align*}
The combined effect $C_G$ and platform-only effect $A_G$ coincide before the concession arrives, under no anticipation.
They need not coincide afterwards. Tsai's first stage uses controls not yet exposed to either event, with baseline
$b=\min(g^1,g^2)-1$. Under an \emph{unconditional} parallel untreated-change restriction for those groups,
\[
 C_G(t)=E[Y_t-Y_b\mid G]-E[Y_t-Y_b\mid G^1>t,G^2>t].
\]
With conditional parallel trends, one must instead adjust and standardize over the target covariate distribution; a raw
unconditional difference is not justified merely by a conditional assumption. Positivity must hold for the selected joint
cohort and comparison observations. The paper develops outcome-regression, inverse-probability and doubly robust moments.
Its sampling setup permits serial dependence within a unit but assumes independent sampling across units, a substantive
limitation for common banking shocks.~\citep{exec28_panel_tsai}

The source's homogeneous-effects decomposition explains why ignoring the concession can matter. If the platform effect is
$\tau_1$, the concession effect is $\tau_2$, and the joint interaction is $\kappa$, the ordinary target-event contrast,
under the additional untreated-change condition needed for that contrast, has the form
\[
 \tau_1+\tau_2(p_T-p_C)+\kappa q_T.
\]
Here $p_T,p_C$ are the fractions newly exposed to the concession between baseline and outcome in the treated and control
groups; $q_T$ is the fraction of the treated group exposed at the outcome date. For synthetic values $\tau_1=40$,
$\tau_2=-30$, $p_T=3/4$ and $p_C=1/4$, the ordinary comparison is $25$, not $40$. If $\kappa=20$ and $q_T=3/4$,
the contrast returns to $40$ by cancellation. Agreement with a favored answer therefore does not diagnose correct
specification. This is a homogeneous illustrative submodel, not identification of unrestricted interactions.
\citep{exec28_panel_tsai,exec28_panel_reference}

Removing contaminated controls identifies the combined effect; it does not automatically separate the interventions.
For $g^1<g^2$ and $t\geq g^2$, a derivation based on Tsai's Appendix A.4 uses the platform effect immediately before the
second event and its subsequent change among units with the same platform cohort but a later concession:
\[
 A_G(t)=C_G(b)+\sum_{s>t}\pi_s\{C_{(g^1,s)}(t)-C_{(g^1,s)}(b)\},\qquad b=g^2-1.
\]
The weights $\pi_s$ are the eligible later-concession cohort shares within the same platform cohort and sum to one.
No anticipation identifies the anchor and donor platform effects; parallel platform-effect changes across concession
cohorts justify transporting the increment. If the anchor is $40$ and the donor effect rises from $45$ to $55$, the
target is $50$. A current combined effect of $85$ must not replace the anchor: doing so gives $95$.
This comparison follows the appendix's explicit anchor and is independently derived here.~\citep{exec28_panel_tsai,exec28_panel_reference}

When $g^2<g^1$, set $b=g^1-1$ and use units with the same concession cohort but a later platform date. Under additive
effects and parallel concession-effect changes across platform cohorts,
\[
 A_G(t)=C_G(t)-C_G(b)-\sum_{s>t}\rho_s\{C_{(s,g^2)}(t)-C_{(s,g^2)}(b)\}.
\]
For synthetic combined effects $95$ and $20$ and a donor concession change from $20$ to $55$, this yields $40$.
The nuisance increment cancels. If concession trajectories differ between the selected cohorts, that cancellation fails.
If adoption dates coincide, this design lacks a period displaying just one effect and does not identify the separate
platform effect. A combined-policy estimand may still be useful, but must be named accordingly.~\citep{exec28_panel_tsai,exec28_panel_reference}

Source precision matters before porting. The selected v3 has inconsistent ordering language near its second stage, and
equation (11)'s target-period entry differs from the pre-concession anchor in Appendix equation (17).
The formulas here follow the appendix's explicit argument and event order; they do not silently declare every printed
expression interchangeable. Section~\ref{sec:exec28-panel-delegates} subsequently inspects the actual
\texttt{fastdid} multiple-event weight path and finds the pre-confounding anchor used in this derivation.
That resolves the selected branch's anchor, while a later-edition comparison, the broader theorem and native
implementation equivalence remain separate questions. The paper names \texttt{fastdid}; pinned repository metadata
independently establishes the author and multiple-event scope. Speed and additional-event claims remain documentation
claims here, not measured performance or a proof for unrestricted treatment histories.~\citep{exec28_panel_tsai,exec28_panel_software}

Inference must retain the full two-stage transformation. For fixed aggregation weights $w$ and first-stage covariance
$\Sigma$, the aggregate variance is $w^{\mathsf T}\Sigma w$, not a sum that discards overlap. Estimated cohort weights
add their own first-order contribution. Tsai invokes smooth parametric nuisance models, linear representations, finite
moments and a multiplier bootstrap. These are conditions to examine, not an automatic justification for dependent branches,
arbitrary machine-learning nuisance fits or rare joint cohorts. The adapter should expose first-stage cells, both event
dates, donor eligibility, aggregation weights and influence functions so the meaning survives serialization.~\citep{exec28_panel_tsai}

\subsection{Balancing an event window does not itself balance stack weights}

Let $a$ index an adoption cohort and $\mathcal A$ retain cohorts for a prespecified event window with adequate clean
controls. Let $n_a^T,n_a^C$ count treated and control units in each sub-experiment, including repeated control appearances
in different stacks. The trimmed treated-share target is
\[
 \theta_e=\sum_{a\in\mathcal A}p_a^T\tau_{a,e},\qquad p_a^T=\frac{n_a^T}{\sum_j n_j^T}.
\]
Keeping $\mathcal A$ and its weights fixed across event time prevents changing composition from masquerading as dynamics.
Constant cohort effects $20$ and $60$ average to $40$ when both are equally represented, but to $20$ when only the early
cohort remains. No effect has decayed. Trimming changes the target population and should be reported as a design choice.
\citep{exec28_panel_wing,exec28_panel_reference}

Wing, Freedman and Hollingsworth show a second problem in the basic pooled stacked comparison. Within each stack, let
the untreated change be $d_a$ in both arms and define $p_a^C=n_a^C/\sum_jn_j^C$. The pooled change contrast is
\[
 \sum_a p_a^T(\tau_{a,e}+d_a)-\sum_a p_a^C d_a
 =\theta_e+\sum_a(p_a^T-p_a^C)d_a.
\]
Parallel trends within every sub-experiment do not cancel the last term when the arms average sub-experiments differently.
This concerns the specified pooled construction; it does not equate every regression described as stacked DiD.
\citep{exec28_panel_wing}

In a synthetic branch example, two stacks contain treated counts $(10,30)$, control counts $(90,10)$, effects $(20,60)$,
and common within-stack untreated changes $(0,100)$. Treated shares are $(1/4,3/4)$ and control shares $(9/10,1/10)$.
The intended effect is $50$. The treated mean change is $125$ and the control change is $10$, so the unweighted pooled
estimate is $115$. Its $65$ excess comes entirely from mismatched trend weights. Even zero causal effects would leave
a false $65$ signal.~\citep{exec28_panel_reference}

Corrective sample weights are one for treated rows and $q_a=p_a^T/p_a^C$ for control rows. Because
$\sum_a p_a^Cq_a=1$, weighted controls have the treated cohort shares. Here control weights are $5/18$ and $15/2$,
the corrected control change is $75$, and the estimate is $50$. Exact row-level means agree with the aggregate algebra.
Large ratios expose weak support and precision costs; they cannot invent controls. A future implementation must reject
zero denominators and distinguish group-count, customer-population and other targets.~\citep{exec28_panel_wing,exec28_panel_reference}

Repeated controls also preserve dependence across stacks. Treating copies of the same branch as independent creates
information absent from the original data. The paper discusses original-unit and unit--stack clustering, and reports a
limited Monte Carlo comparison. This does not establish universal coverage for cross-branch shocks or few policy clusters.
Original IDs must survive stacking, and the covariance specification must state which common shocks it permits.
Correcting point-estimate weights alone does not qualify uncertainty.~\citep{exec28_panel_wing}

\subsection{Conditional reuse across the R and Stata routes}

The inspected stacked tutorial contains R and Stata examples, sample-share weights, and regression calls clustered on
the original state identifier. It is a useful prototype template subject to a pinned environment and membership/weight
tests, not a complete adapter. Its R inclusion expression also contains an \texttt{adoptionTime == TRUE} branch alongside
the intended treated, late and missing-date branches. Numeric-date encoding is therefore a specific review obligation.
The Stata weight example sums within-group row indices, proportional to counts only under appropriate balance.
Balanced and unbalanced fixtures are needed before translation. The complete selected tree has no named license file;
that bounded observation is not a legal conclusion about all possible permissions.~\citep{exec28_panel_software}

The nonabsorbing route addresses a different question. Pinned \texttt{nonabsdid} 0.4.1 metadata describes a wrapper around
\texttt{DIDmultiplegtDYN}, \texttt{PanelMatch} and \texttt{fect}, with an optional \texttt{fixest} reference.
A shared plotting schema is useful, but does not equate targets, histories or uncertainty procedures. Consider adopting
output conventions while inspecting every delegate and retaining method-specific fields. Direct adoption of the comparison
as a common estimand is premature. The metadata declares MIT plus a file identifying the copyright holder; dependency
rights and complete behavior remain separate checks.~\citep{exec28_panel_software}

Pinned \texttt{pampe} 1.1.2 documentation describes a single affected unit, donor outcome matrix, subset regression selection
and counterfactual prediction, with control and time placebos. Its documented defaults use AICc and a \texttt{leaps}
dependency. This suggests a comparator for a regional intervention with a credible stable donor relation, not a cohort ATT
obtained by renaming output. The manual's model-size discussion and actual delegate need reconciliation; the original
panel-approach argument and package paper require separate inspection. Metadata declares GPL-2. An isolated callable
comparison is a proposal option, not an approved port or redistribution decision.~\citep{exec28_panel_software}

These choices imply a concrete CLI/MCP contract: original unit IDs, outcome and event clocks, competing interventions,
reversal policy, baseline, clean-control and cohort rules, aggregation population, covariance assumptions and package identity.
A result must expose its supported target and excluded cells as well as estimates. An LLM can help elicit intervention
calendars and explain why a proposed control is already exposed; it cannot infer unobserved parallel trends or certify a
causal diagram from a plot. No acquired statistical engine or sibling repository was executed for these comparisons.

The source boundary is explicit. Roth is exact arXiv v2 submitted 16 January 2026 with a 21 January cover; Tsai is exact
v3 submitted 21 January 2025 with a 22 January cover. The stacked original is the retained January 2024 NBER working paper
from a mutable endpoint; its cover does not prove historical byte identity. The cutoff remains 7 October 2026.
Later-edition clarification, remaining originals, delegate inspection, dependency rights and rendered visual review remain
local obligations. The reviewed panel additions do not close the overall software or proposal program.
